\documentclass{article}
\usepackage[T1]{fontenc}
\usepackage{lmodern}
\usepackage{graphicx}
\usepackage{amsmath,amssymb}
\usepackage[a4paper,margin=2cm]{geometry}
\usepackage[absolute,overlay]{textpos}
\setlength{\TPHorizModule}{1mm}
\setlength{\TPVertModule}{1mm}
\usepackage[indonesian]{babel}
\usepackage{hyperref}
\setlength{\emergencystretch}{2em}
% unit: Halaman 60

\begin{document}

% === Halaman 60 (llm) ===

\section*{Kode Program RC4 (dalam Bahasa Python)}

\begin{verbatim}
from Crypto.Cipher import ARC4
import base64

def RC4_encrypt(plainteks : str, kunci : bytes) -> str:
    cipher = ARC4.new(kunci)
    cipherteks = cipher.encrypt(plainteks.encode())
    return base64.b64encode(cipherteks).decode()

def RC4_decrypt(cipherteks : str, kunci : bytes) -> str:
    cipher = ARC4.new(kunci)
    cipherteks_decoded = base64.b64decode(cipherteks)
    return cipher.decrypt(cipherteks_decoded)

pesan = input("Ketikkan pesan rahasia:")
kunci = input("Ketikkan kata kunci: ")
kunci = bytearray(kunci, 'utf-8')
ciphertext = RC4_encrypt(pesan, kunci)
print(f"Cipherteks: {ciphertext}")
plaintext = RC4_decrypt(ciphertext, kunci)
print(f"Plainteks semula: {plaintext}")
\end{verbatim}

\end{document}
